\documentclass[runningheads]{llncs}
\usepackage[T1]{fontenc}
\usepackage{graphicx}
\usepackage{booktabs}
\usepackage{amsmath}
\usepackage{url}
\usepackage{xcolor}
\usepackage{array}
\renewcommand\UrlFont{\color{blue}\rmfamily}
\newcommand{\adapter}{a GRPO LoRA adapter (identifier withheld for review)}

\begin{document}
\sloppy

\title{Receipt-Bound Held-Out Evaluation of a GRPO-Trained Language Model Across Fourteen Agentic and Reasoning Suites}
\titlerunning{Receipt-Bound Evaluation of a GRPO-Trained LLM}

\author{Anonymous Author(s)}
\authorrunning{Anonymous}
\institute{Anonymous Institution (withheld for double-blind review)}

\maketitle

\begin{abstract}
Benchmark results for reinforcement-learning post-trained language models are
usually reported as single headline numbers. The report rarely says how much of
each suite was actually graded, whether the original benchmark or a substitute
was run, or whether the number can still be traced to the run that produced it.
We report a held-out evaluation of one Group Relative Policy Optimization
(GRPO)-trained actor against fourteen agentic and reasoning benchmark suites
under a receipt-bound, fail-closed governance regime. The actor is fixed across
every lane, each suite is graded only by its own native evaluator at a pinned
revision, and every reported figure is bound to a write-once receipt that a
deterministic ledger check recomputes. Four scopes closed complete: LAB-Bench
public split (1,967/1,967 items, accuracy 0.2288), AgentDojo benign utility
(97/97 episodes, utility 0.9072), VerilogEval (312/312 problems, 41.35\%
pass@1) and Omni-MATH (4,428/4,428 dispositions, official accuracy 51.31\%
reproduced). SWE-bench Pro carries a full original-contract score of
2/731 = 0.274\% pass@1 with every non-native outcome retained in the
denominator. The remaining lanes are reported at their true terminal states:
partial with coverage, blocked on cloud quota, or closed as externally blocked
with a named reopen condition. A separately authorised base-model rerun of
BankerToolBench scored 0.0 over 100/100 trials, which an audit traces to
tool-dialogue collapse rather than domain reasoning. Eleven of eleven
deterministic ledger checks pass. Original-contract and replacement-scope
numbers are never pooled, no cross-suite aggregate is computed, and no
improvement over any baseline is claimed.
\keywords{Large language models \and Reinforcement learning \and Benchmark evaluation \and Agentic evaluation \and Reproducibility \and Evaluation governance}
\end{abstract}

\section{Introduction}

Reinforcement learning with verifiable rewards has become a standard stage in
post-training large language models (LLMs). Group Relative Policy Optimization
(GRPO)~\cite{shao2024deepseekmath,guo2025deepseekr1} is a common choice because
it removes the value critic and derives each completion's advantage from the
reward statistics of its sampling group. Once a model has been trained this way,
the question most readers ask is how well it does on held-out benchmarks. The
answer is usually a table of headline numbers.

Such tables hide three things. First, many recent agentic and reasoning suites
are \emph{multi-scope}: a public split and a private original contract, or a
benign-utility scope and an adversarial scope. A number from one scope is easily
read as a number for the suite. Second, many suites are quota-limited, gated
behind private payloads, or expensive, so a run may grade only a fraction of the
suite. Reports tend either to drop such suites or to quote a score over whatever
subset graded, which hides the coverage. Third, the chain from a reported number
back to the run that produced it is often broken: the sampler path is gone, the
execution script was deleted, or the arithmetic was done by hand. Earlier work
showed how sensitive reinforcement-learning results are to seeds and
implementation details~\cite{henderson2018deep}, and reproducibility
checklists~\cite{pineau2020improving}, model cards~\cite{mitchell2019modelcards}
and datasheets~\cite{gebru2021datasheets} set norms for disclosure. The
evaluation layer of RL post-trained LLMs has not yet been held to the same
standard.

This paper reports one campaign built to that standard. A single GRPO-trained
actor is evaluated against fourteen suites (lanes E1--E14) spanning software
engineering, tool-using agents, web navigation, binary analysis, scientific
question answering, machine-learning engineering, agent safety and utility,
hardware design, application deployment, sequential games and competition
mathematics. The contributions are:
\begin{itemize}
\item a \textbf{one-actor, native-evaluator design}: no per-lane fine-tuning,
prompt tuning or evaluator substitution, so differences between lanes reflect
the tasks rather than per-lane preparation;
\item a \textbf{receipt-bound, fail-closed governance regime} in which every
figure traces to a write-once receipt, preparation, execution and scoring are
tracked separately, and a lane that cannot meet its preconditions records that
state instead of substituting a neighbouring benchmark;
\item \textbf{per-lane results at their true terminal states}, including
complete scopes, full-suite original-contract scores, partial lanes with their
denominators, and externally blocked lanes with reopen conditions;
\item an \textbf{explicit account of threats} that governance could not remove,
including a sampler path that no longer exists and execution sources lost after
the fact.
\end{itemize}
The companion training study, which examines GRPO's zero-variance
signal-starvation regime, is reported separately and is referred to here only
for context.

\section{Campaign Design}

\subsection{The Evaluated Actor}

One actor is evaluated in every lane. The base model is
\texttt{Qwen/Qwen3.6-35B-A3B}, a sparse mixture-of-experts model, at revision
\texttt{995ad96e}. The trained parameters are \adapter, produced by GRPO
training on the Tinker platform~\cite{thinkingmachines2024tinker} with low-rank
adaptation~\cite{hu2022lora}, and served in bfloat16. No lane received a
separately fine-tuned checkpoint, per-suite prompt engineering beyond what the
suite's own harness prescribes, or tuning against its evaluation set.
Two consequences follow. Differences between lanes reflect one fixed policy
meeting fourteen different interfaces, observation spaces, tool conventions
and graders, not what a specialised actor could achieve. And the campaign
cannot say anything about seed variance: where a lane sampled more than one
completion per task, it did so because the suite's protocol requires it.

Inference conditions are fixed per lane and recorded in that lane's receipt
rather than imposed globally, because the sampling regime is part of the
result. For example, the AgentDojo lane ran at temperature 0 with a
32,768-token context and 4,096 output tokens per request, while the VerilogEval
lane sampled at temperature 0.2, top-$p$ 0.95, 1,024 maximum output tokens and
seed 1816, one sample per problem.

\subsection{Lanes, Contracts and Replacement Scopes}

Each lane is defined by one original benchmark contract and, where the
original could not be obtained, one named replacement scope
(Table~\ref{tab:lanes}). The replacement is always named and reported
separately. It is never presented under the original suite's name, and the two
are never pooled, because a replacement is a different suite with a different
interface and often a different grader.

Each lane is graded only by its suite's native evaluator at a pinned revision.
No lane replaces the upstream grader with a re-implemented metric, because a
substituted metric would no longer be evidence about the benchmark it claims to
measure.

\subsection{Governance: Receipts, Budgets and Decontamination}

Every graded run follows the pipeline in Fig.~\ref{fig:pipeline}. A sealed,
hash-pinned request authorises execution under a bound budget. The suite's
native grader produces rows and its own denominator. A write-once JSON receipt
records the actor identity, sampler path, suite revision, inference settings,
spend and decontamination status. A deterministic ledger check then recomputes
the arithmetic before any figure enters a results table. If any link is
missing, the chain stops and the lane reports no number.

\begin{figure}[t]
\centering
\includegraphics[width=\textwidth]{figures/fig_receipt_pipeline_cropped.pdf}
\caption{Receipt-bound evaluation chain (top). A figure is reported only if
every link exists. When a link is missing, e.g.\ a lost execution source
(bottom), the chain stops and the lane reports no number.}
\label{fig:pipeline}
\end{figure}

Three quantities are tracked separately: what was \emph{prepared} (suite
inventoried, harness built, manifest validated), what was \emph{executed}
(tasks actually run against the actor) and what was \emph{scored} (tasks with a
native, verifiable verdict). A lane with heavy preparation and little scoring
is a different situation from a lane that was never attempted, and the records
keep them apart.

Each paid run is authorised with a per-run spend cap and timeout. A later
authorisation removed the cumulative campaign cap but kept every per-run
resource and timeout bound, so spend is gated at each launch by a preflight
check rather than by a campaign-wide ceiling. Receipts also carry a
decontamination status. For example, the AgentDojo receipt records
\texttt{TRAINING\_INVENTORY\_ABSENT} and states its claim boundary verbatim as
benign task utility only, with no prompt-injection security or held-out claim.

The campaign is \emph{fail-closed}. When a precondition cannot be met, e.g.\ a
withheld credential, a denied cloud quota or an unreleased private bundle, the
lane records that state. Every lane ends in a named terminal state:
\texttt{COMPLETE}, \texttt{CLOSED\_PARTIAL},
\texttt{REBUILD\_READY\_LAUNCH\_PENDING},
\texttt{AMENDMENT\_ACCEPTED\_LAUNCH\_PENDING}, \texttt{PENDING\_QUOTA}, or
\texttt{CLOSED\_EXTERNAL} with a reopen condition.

Two reporting rules hold throughout. Original-contract and replacement-scope
numbers are never pooled. And no cross-suite mean, median or composite index is
computed, because native denominators differ in size, in what they count and in
what they exclude, so an aggregate over them would have no referent.

\section{Results}

Table~\ref{tab:lanes} gives every lane at its terminal state, and
Fig.~\ref{fig:lanes} summarises the states. Figures are quoted with their own
denominators and coverage; rows are not comparable with one another.

\begin{table}[t]
\caption{All fourteen lanes. O = original contract, R = replacement scope.
Coverage and score use each suite's native denominator. ``---'' means no
score exists. Rows are never pooled or averaged.}
\label{tab:lanes}
\centering
\scriptsize
\setlength{\tabcolsep}{3pt}
\resizebox{\textwidth}{!}{%
\begin{tabular}{@{}l l c l l l@{}}
\toprule
Lane & Suite & & Coverage & Score & Terminal state \\
\midrule
E1  & SWE-bench Pro~\cite{deng2025swebenchpro} & O & 731/731 & 2/731 = 0.274\% & full suite \\
    & SWE-bench Multilingual & R & 35/300 graded & --- & rebuild ready, launch pending \\
E2  & FrontierSWE & O & 1/17 & 0.8628 (1 task) & partial \\
    & CORE-Bench & R & 0 graded (45 setups) & --- & launch pending \\
E3  & SDAB (private bundle) & O & --- & --- & closed, external \\
E4  & BankerToolBench & O & 1/100 & 0.3115 (recovery) & closed, partial \\
E5  & APEX-Agents & O & 7/480 & --- (not suite score) & partial \\
    & Tau3 & R & 20/97 (20.62\%) & --- & rebuild ready, launch pending \\
E6  & WebArena & R & 0/812 & --- & pending quota \\
E7  & BinaryAudit & O & 1/46 attempted & --- (agent error) & closed, external \\
E8  & LAB-Bench public~\cite{laurent2024labbench} & R & 1,967/1,967 & 0.2288 & complete \\
E9  & MLE-bench & O & 40/75 (53.33\%) & null & partial \\
    & MLDevBench & R & 0/34 & --- & pending quota \\
E10 & AgentDojo benign~\cite{debenedetti2024agentdojo} & R & 97/97 & 0.9072 & complete \\
E11 & VerilogEval~\cite{liu2023verilogeval} & O & 312/312 & 129/312 = 41.35\% & complete \\
E12 & AppBench & O & --- & --- & closed, external \\
E13 & BALROG & R & 13/255 & --- & launch pending \\
E14 & Omni-MATH~\cite{gao2024omnimath} & R & 4,428/4,428 & 2,271/4,428 = 51.31\% & terminal-complete \\
\bottomrule
\end{tabular}}
\end{table}

\begin{figure}[t]
\centering
\includegraphics[width=\textwidth]{figures/fig_lane_status_cropped.pdf}
\caption{Terminal state of every lane: complete, partial (awaiting launch or
closed with partial evidence), or blocked (external closure or cloud quota).}
\label{fig:lanes}
\end{figure}

\subsection{Complete Scopes}

\textbf{E8, LAB-Bench (public split).} All 1,967 of 1,967 items across eight
categories were evaluated, with 450 correct: accuracy
$450/1967 = 0.2288$. Per-category accuracy ranges from 0.7869 on TableQA
(192/244) to 0.0 on CloningScenarios (0/33). One receipt field matters for how
this figure is read. For 1,259 of the 1,967 responses the native pipeline
extracted no answer, and these are counted as incorrect. The figure is
therefore an end-to-end harness accuracy, not accuracy over the parseable
subset.

\textbf{E10, AgentDojo (benign-utility scope).} Against AgentDojo v1.2.2, all 97
episodes completed and 88 passed their utility predicates, for a benign utility
of $88/97 = 0.9072$ (banking 16, slack 21, travel 20, workspace 40 episodes). The
list of missing evaluation IDs is empty and no native errors were handled.
Completion and utility are different quantities, so the lane is ``complete in
coverage, 0.9072 in benign utility'', never ``97/97 correct'', and it says
nothing about security. An earlier attempt stopped at 94/97 because of a
collector fault in trace-path comparison, not a model failure. The repair
touched only that path rewrite, passed 18 offline tests, and a continuation run
closed the remaining three episodes.

\textbf{E11, VerilogEval.} All 312 pinned problems were evaluated with the
upstream harness, Icarus Verilog and Verilator: 129 passed, or 41.35\% pass@1.
Of these, 67/156 were code-completion problems and 62/156
specification-to-RTL problems. A reference sweep showed that one problem cannot
be passed by any implementation because its test bench fails to elaborate.
Excluding it gives 129/311, but no benchmark-level justification for the
exclusion was recorded, so 129/312 remains the canonical figure and 129/311 is
kept only as a non-canonical sensitivity. The 150 responses from which no code
could be extracted count as failures in the denominator.

\textbf{E14, Omni-MATH (replacement for FrontierMath).} The native scorer
accepted 4,426 of 4,428 rows and reproduced an official accuracy of
$2{,}271/4{,}428 = 51.31\%$. For the two rejected rows, the judge's output was
truncated before any equivalence verdict, so omitting them is correct official
behaviour. Both already count as not correct in the denominator, so the score
impact is nil. The lane closes terminal-complete with all 4,428 dispositions
recorded, and no re-judge was run, since that would depart from the official
scorer.

\subsection{Original-Contract Results}

\textbf{E1, SWE-bench Pro.} The actor resolved 2 of 731 tasks in the full test
split, a pass@1 of 0.274\%. Within that denominator are 713 native evaluations,
14 generation failures and 4 lost generation artifacts. These 18 non-native
outcomes stay in the denominator, so the figure is a rate over all 731 tasks.
The lane records one fidelity caveat: the gold-resolve timeout was 1,200\,s
rather than the official 120\,s, because of an emulation constraint in the local
execution environment.

The other original-contract figures are partial and none is a suite score.
FrontierSWE (E2) evaluated 1 of 17 tasks, with a replay-normalised score of
0.8628. BankerToolBench (E4) evaluated 1 of 100 tasks, with a verifier recovery
metric of 0.3115 on that single task. APEX-Agents (E5) natively scored 7 of 480
tasks, averaging 0.079365. Its prefix mean of 0.050505 is an attempt summary
that counts unscored attempts as zero, not a 480-task score. BinaryAudit (E7)
records 1 of 46 attempts, with a verifier reward of 0.0 after an agent error and
no grade. MLE-bench (E9) has native grades for 40 of 75 competitions (53.33\%
coverage) but a null suite score, because coverage without complete grading is
not a score.

\subsection{Partial and Quota-Blocked Replacement Scopes}

Four replacement scopes are staged, partly measured, and waiting on a launch
step. SWE-bench Multilingual (E1) has 35/300 graded. CORE-Bench (E2) has all 45
capsule setups done and 0 graded; its driver passes 12/12 offline tests and it
is blocked on an owner-level cloud IAM binding. Tau3 (E5) has 20/97 tasks
cleaned (20.62\%). BALROG (E13) has 13/255 episodes recorded. Two scopes are
blocked only on cloud quota. WebArena (E6) is at 0/812 with its 812-task
inventory verified, blocked on a 1-of-16 vCPU regional quota. MLDevBench (E9)
is at 0/34 with an incomplete runtime image, blocked on a 1-of-8 vCPU quota.
Both were rechecked live and recorded NO-GO.

\subsection{Externally Blocked Lanes}

Lanes whose original suite depends on a private or hosted resource were closed
as \texttt{CLOSED\_EXTERNAL}: SDAB (E3, private bundle), BinaryAudit (E7,
private payload), LifeSciBench (E8-original), AgentHarm (E10-original, private
tasks), AppBench (E12, deployment), OpenReward Games (E13-original) and
FrontierMath (E14-original, hosted evaluation). Each has a closure record and
a reopen condition that depends on a provider response. Access requests
were made through public issue trackers, dataset discussions and email, and
are reported with their delivery status. Where email delivery could not be
verified, the record says so.

\subsection{A Diagnostic Zero: the BankerToolBench Rerun}

A separately authorised rerun of BankerToolBench used the base model (a
zero-initialised LoRA snapshot, equivalent to base sampling). It completed
100/100 trials in 3 hours 29 minutes and cost \$15.07 against a \$55.91 cap.
Mean reward was exactly 0.0, with all 100 trials in the 0.0 bucket. An audit of
every trajectory found the mechanism. Agents stop after 2 to 14 steps (median
about 6) without producing a deliverables directory, so the verifier returns 0.0
with ``No deliverables found''. At least 27 trajectories end in role-token
degeneration of the form \texttt{user user assistant \dots}. The zero therefore
measures \emph{tool-dialogue collapse}, not finance reasoning. The harness
fails closed and gives a clear reason for every zero instead of manufacturing
partial credit. It also means no BankerToolBench suite score is currently
obtainable for this actor family. The rerun is not set against the earlier
single-task E4 figure, since the two were run under different configurations.

\subsection{Ledger Verification}

The ledger's arithmetic and provenance were re-verified by eleven deterministic
code checks, and all eleven pass with zero failures. They cover exact divisions
(e.g.\ $67 + 62 = 129$ for E11's two framings), receipt presence,
sealed-artifact hash equality, the claim that exactly E8, E10 and E11 are the
strictly complete replacement scopes, and absence probes for the lost execution
sources. A reviewer can therefore verify every number's arithmetic and
provenance without the compute that produced it. The same deterministic rule
applies to model-assisted triage used during the campaign: arithmetic and
string lookups stay in code, and model judgment is used only for
classification.

\section{Threats to Validity}

\textbf{No base-model comparator.} No lane has a matched base-model run, so no
campaign figure is evidence of what the GRPO step contributed. The one
base-model result (E4 rerun) is a diagnostic, not a paired baseline.

\textbf{Native denominators are not comparable.} 41.35\% on 312 problems and
0.274\% on 731 tasks measure different capabilities under different scoring
rules. No average, ranking or composite is offered, and the rule applies
within lanes too.

\textbf{A sampling path that no longer exists.} The VerilogEval scores were
produced by sampling the adapter through the training platform's managed
sampler. That sampler checkpoint was later purged, so a rerun would require
serving the published mirror of the same weights through a different sampler.
The result stays receipt-bound, but its sampling path cannot be re-run, and a
rerun would compare the same weights through a different sampler, a divergence
that has not been measured.

\textbf{Lost execution sources.} A working-directory deletion removed the
execution sources of several late recovery launches, while their sealed
requests, validation receipts and test logs survive with matching hashes. Under
a lost-source policy, one launcher was faithfully re-implemented in a tracked
path and passes 10/10 targeted regressions offline. A sealed request without
its execution source is not a rerun, and re-implementation cannot restore what
was lost. From then on, all execution sources are kept in tracked paths.

\textbf{Scope and harness fidelity.} The benign AgentDojo scope measures no
adversarial robustness. SWE-bench Pro ran with a longer gold-resolve timeout
than the official setting. LAB-Bench accuracy counts unparsed responses as
incorrect. Each caveat is reported alongside its number rather than corrected
after the fact.

\section{Conclusion}

A governance-first campaign produced complete, receipt-bound results on four
scopes, a full original-contract score on SWE-bench Pro, and an auditable
terminal state for every remaining lane, each with a named gate and reopen
condition. Its value lies in the rules it enforced. Original-contract and
replacement-scope numbers are never pooled, no cross-suite aggregate is
computed, and no improvement over a baseline is claimed. The strongest complete
positive figure, 41.35\% pass@1 on VerilogEval at full coverage, is a
measurement of one actor on one suite, with provenance recorded down to its
sampler path. We suggest that evaluations of RL post-trained LLMs report
prepared, executed and scored coverage separately, bind each figure to a
receipt, and record blocked lanes as blocked rather than dropping them. Future
work will add matched base-model runs, so that per-suite post-training effects
can be estimated, and will complete the staged lanes once their gates clear.

\bibliographystyle{splncs04}
\bibliography{refs}

\end{document}
